\subsection{Deterministic {\em vs} Stochastic Dynamics:\\ A theory of island biogeography in the limit of $S \rightarrow \infty$}
As you know, in 1967, MacArthur \& Wilson published a book ``The Theory of Island Biogeography'' (TIB) to explain the factors controlling species diversity on {\em islands} at certain distance from the continent, the {\em mainland}. The stochastic model underlying this theory will be reviewed around the 7th week. A simpler version of it makes the following assumptions:
\begin{enumerate}
  \item {\bf $S \rightarrow \infty$:} The number of species, $S$, in the mainland is very large. Technically, $S$ has tended to $\infty$.
  \item {\bf Species equivalence:} Species arrive in an island at the same rate, $\lambda$, the colonization rate, and, they get extinct at the same rate, $\mu$, as well.
  \item {\bf Species independence:} Species arrivals and extinctions are random independent events. 
\end{enumerate}

Given the previous assumptions, address the following questions:
\begin{enumerate}
\item Suppose you have an island with a certain number of species $n$. What is the deterministic equation governing the temporal evolution of this number? Use this equation to calculate the number of species, $n^{\star}$, at the colonization-extinction equilibrium? 

\item The probability a new species arrives in time $dt$ is given by $\Pb[n \rightarrow n+1] = \lambda\ dt$. Since all species become extinct at the same rate, $\mu$, independently, the probability of loosing any of them to extinction is given by $\Pb[n \rightarrow n -1] = \mu\ n\ dt$. Use this probabilities to build the set of differential equations for the temporal evolution of the probability distribution describing $n$. This equation is the stochastic counterpart of the previous deterministic equation. 

\item Calculate the stationary distribution for the number of species at equilibrium, $\Pb^{\star}(n)$. What is the expected value for this number?

\item Consider a system of $K$ islands undergoing the colonization-extinction process described by the three assumptions above.  Assume there is no inter-island colonization so that all islands receive species directly from the mainland at rate $\lambda$. Assume also the system has reached equilibrium. If species on an island are detected with probability $a$, what is the probability of observing $n_i$ species on the $i$-th island? 

\item Assume you visited the archipelago of $K$ islands for a survey and observed $n_1, \dots, n_K$ species on the 1st to the $K$th islands, respectively. Assume again that there is no inter-island colonization so that all islands receive species directly from the mainland at rate $\lambda$. The sample of species observed with detectability $a$ across islands can be represented by the vector $\vec{n} = (n_1, \dots, n_K)$. Assume the process underlying this sample is described by the three assumptions above and that the whole archipelago is at colonization-extinction equilibrium. Calculate the probability of obtaining the sample $\vec{n}$, $\Pb^{\star}(\vec{n}) | \lambda, \mu)$. 

\item The probability $\Pb^{\star}(\vec{n}) | \lambda, \mu)$ is sometimes called a sampling probability distribution. It can be regarded also as the likelihood of obtaining the sample $\vec{n}$ given that model parameters are $\lambda$, $\mu$, and $a$. If we define 
\begin{equation}
  F(a, \lambda, \mu \,|\, \vec{n}) \equiv \Pb^{\star}(\vec{n} \,|\, a, \lambda, \mu) 
\end{equation}
and take logarithms, we obtain: 
\begin{equation}
  \log \left( F(a, \lambda, \mu \,|\, \vec{n} ) \right) = \log \left( \Pb^{\star}(\vec{n} \,|\, a, \lambda, \mu) \right) 
% * <dalonso@ceab.csic.es> 2017-10-17T15:31:45.487Z:
%
% ^.
\end{equation}
which is called the loglikelihood function. By defining the auxiliary variable $x \equiv a \lambda/\mu$, calculate the parameter $x^{\star}$ that maximizes the probability of obtaining the sample vector $\vec{n}$. 
\end{enumerate}

